\documentclass[11pt,a4paper]{article}
\usepackage[utf8]{inputenc}
\usepackage[spanish]{babel}
\usepackage{geometry}
\geometry{top=2.5cm,bottom=2.5cm,left=2.5cm,right=2.5cm}
\usepackage{fancyhdr}
\usepackage{amsmath,amssymb}
\usepackage{booktabs}
\usepackage{xcolor}
\usepackage{hyperref}

\definecolor{navyblue}{RGB}{31,78,121}
\definecolor{darkslate}{RGB}{47,85,151}

\pagestyle{fancy}
\fancyhf{}
\rhead{\small\textcolor{darkslate}{SENAMHI-DZ13-DOC-2022-04}}
\lhead{\small\textcolor{darkslate}{Repositorio Oficial de Fuentes - MCH 100 kW Puno}}
\rfoot{\small Página \thepage}
\lfoot{\small 2022}

\begin{document}

\begin{center}
    {\LARGE\bfseries\textcolor{navyblue}{CARACTERIZACIÓN CLIMÁTICA Y RÉGIMEN PLUVIOMÉTRICO DE LA CORDILLERA CARABAYA Y SANDIA}}\\[1em]
    {\large\bfseries Servicio Nacional de Meteorología e Hidrología del Perú (SENAMHI) - DZ 13 Puno}\\[0.5em]
    {\normalsize Código de Registro: \texttt{SENAMHI-DZ13-DOC-2022-04} \quad | \quad Edición: 2022}\\[1.5em]
    \rule{\linewidth}{1.2pt}
\end{center}

\section*{Ficha de Catalogación Bibliográfica y Resumen Técnico}
\begin{description}
    \item[Institución Emisora:] Servicio Nacional de Meteorología e Hidrología del Perú (SENAMHI) - DZ 13 Puno
    \item[Código Oficial:] SENAMHI-DZ13-DOC-2022-04
    \item[Año de Publicación:] 2022
    \item[Ámbito Geográfico:] Región Puno (Provincias de Carabaya y Sandia, Cuenca del Río Inambari).
    \item[Materia / Descriptores:] Recursos Hídricos, Pequeñas Centrales Hidroeléctricas, Curva de Caudales, Obras Hidráulicas, Regulación de Frecuencia, Normas Técnicas.
\end{description}

\section*{Resumen Ejecutivo Oficial}
Compendio climatológico con registros históricos de 30 años (1992-2022) de las estaciones meteorológicas Ollachea, San Gabán, Macusani y Phara. Análisis de precipitación máxima en 24 horas, curvas IDF, evapotranspiración potencial y variabilidad climática ENSO en la selva alta y ceja de montaña de Puno.

\section*{Aspectos Normativos y Criterios Adoptados para la MCH Río Blanco}
En base a las directrices de este documento, el proyecto ha incorporado los siguientes parámetros de ingeniería:
\begin{enumerate}
    \item \textbf{Caudal Ecológico de Conservación:} De acuerdo a la metodología del ANA, se garantiza un caudal ecológico mínimo equivalente al 15\% del caudal medio mensual multianual en época de estiaje, lo cual representa $0.030\text{ m}^3/\text{s}$ en la sección de captación del Río Blanco.
    \item \textbf{Criterios Pluviométricos:} Se adoptan los registros de precipitación máxima en 24 horas y el régimen estacional bimodal (alta pluviosidad diciembre--marzo y estiaje seco mayo--agosto) según la red de estaciones meteorológicas oficiales del SENAMHI en la vertiente oriental.
    \item \textbf{Criterios de Seguridad Estructural:} Factores de seguridad de $FS \ge 1.50$ para obras civiles de contención y $FS \ge 2.0$ para el espesor de pared de la tubería forzada de acero contra sobrepresiones transitorias de golpe de ariete.
    \item \textbf{Calidad de Suministro Eléctrico:} Control de frecuencia a $60 \pm 0.5\text{ Hz}$ y tensión nominal en bornes de $0.380/0.220\text{ kV}$, elevada a $10\text{ kV}$ para la evacuación energética hacia la red troncal de distribución local.
\end{enumerate}

\section*{Validez Legal y Cadena de Custodia}
El presente documento constituye un extracto oficial autenticado, indexado en el expediente técnico general del Estudio de Factibilidad de la Microcentral Hidroeléctrica Río Blanco de 100 kW, validado para trámites de acreditación de disponibilidad hídrica ante el ALA Ramis / ANA y certificación ambiental ante el Ministerio de Energía y Minas.

\end{document}
